\documentclass{beamer}
\input{tbook-deck}
\title{One base, every lift splits}
\subtitle{A negative answer to Kang--Jäger's Question 1.10}
\author{Haobo Ma\inst{1} \and Wenlin Zhang\inst{1,2}}
\institute{\inst{1}Chrono AI, Singapore\and\inst{2}National University of Singapore, Singapore}
\date{Version 2, 2026}
\begin{document}
% Manuscript: Title and abstract
\begin{frame}
\titlepage
\note{The question asks for some local rule preserving minimality of the entire lift. The result constructs one fixed base defeating every length and every bipermutative rule.}
\end{frame}

% Manuscript: Sections 1–4
\begin{frame}{Route}
\tableofcontents
\note{Begin with a local rule and the order of quantifiers. Then construct the blocks and explain the invariant label and evidence.}
\end{frame}

\section{Question}
\subsection{Local rules}

% Manuscript: Section 1, definitions; Section 4, finite-state representation
\begin{frame}{Two inputs need not mean two components}
The adjacent XOR rule:
\[
y_i=x_i+x_{i+1}\pmod 2,\qquad x_{i+1}=x_i+y_i\pmod 2.
\]
\begin{itemize}
\item The output and $x_0$ uniquely reconstruct both directions.
\item Two starting bits give two inputs.
\item Shifting changes the bit at the origin; we need an invariant label.
\end{itemize}
\note{Use the simplest two-bit rule to introduce finite memory. Two preimages above each output do not by themselves imply dynamical splitting.}
\end{frame}

% Manuscript: Section 1; Kang–Jäger Question 1.10 and Theorem 1.9
\begin{frame}{The order of choices matters}
Kang--Jäger's Question 1.10: given a minimal or strictly ergodic base, does some permutative rule of length at least two preserve that property in its entire inverse image?
\medskip
\begin{itemize}
\item Earlier: fix a rule, then choose a base that defeats it.
\item Here: choose one base before every length and rule.
\item The whole inverse image matters, not a selected component.
\end{itemize}
\note{Bipermutative means bijective in either endpoint with the other letters fixed. The range is $\ell\geq2$. Compare the order of choices with Kang--Jäger, Theorem 1.9(b),(c).}
\end{frame}

% Manuscript: Theorem 1
\begin{frame}{The same base for every code}
\begin{bookresult}{thm-main}
There is an aperiodic regular Toeplitz, strictly ergodic $Y$ such that every $\ell\geq2$ and every bipermutative $h$ admit
\[
\Theta_h:\eta_h^{-1}(Y)\simeq Y\times\{0,1\}^{\ell-1},
\]
\[
\operatorname{pr}_1\Theta_h=\eta_h,\qquad
\Theta_h S=(S\times\mathrm{id})\Theta_h.
\]
\end{bookresult}
\note{Minimal means every orbit is dense. Strict ergodicity also requires a unique invariant probability. Toeplitz means each coordinate has its own period. The statement includes a homeomorphism, its first coordinate and conjugacy, not just nonminimality.}
\end{frame}

\section{Construction}
\subsection{Recognizable blocks}

% Manuscript: Section 2, equation (1)
\begin{frame}{Factorial blocks}

\[
U_0=0,\quad V_0=1,\quad m_n=(n+1)!\quad(n\geq0).
\]
\[
U_{n+1}=U_n^{2m_n}V_n^{m_n},
\]
\[
V_{n+1}=U_n^{m_n}V_n^{2m_n}.
\]
\[
(U_1,V_1)=(001,011),\quad H_{n+1}=3m_nH_n.
\]
Rows tiled at every level form $Y=\bigcap_{n\geq0}Y_n$; nested compactness ensures nonemptiness.
\note{Powers denote concatenation. Both words at level $n$ have length $H_n$. Level two has length $18$. Nonemptiness follows from compactness across all levels, not finite checks.}
\end{frame}

% Manuscript: Lemma 2; Section 2, phase identity
\begin{frame}{The symbols reveal the boundaries}

\begin{bookresult}{lem-tiling}
Each marked tiling is unique and continuous; on $Y$ the levels refine compatibly.
\end{bookresult}
\[
0^{2m_n}1^{m_n}\quad\hbox{or}\quad0^{m_n}1^{2m_n}.
\]
A $10$ in the type sequence marks exactly a parent start; run lengths identify its type.
\[
\theta_n(Sy)=\theta_n(y)+1\pmod {H_n}.
\]
$S^py=y$ would force every $H_n\mid p$, excluding positive periods.
\note{Continuity follows because the continuous reading bijection from a compact tiling space to the Hausdorff row space is a homeomorphism. Existing formal sources give injectivity at all levels and continuous decoding on the intersection; the full finite-level inverse statement retains its written proof.}
\end{frame}

% Manuscript: Lemma 3
\begin{frame}{Every local pattern returns}

\begin{bookresult}{lem-minimal}
$Y$ is minimal.
\end{bookresult}
\begin{itemize}
\item For $n\geq1$, every level-$(n+2)$ word contains all type pairs.
\item A window of length at most $H_n$ meets at most two level-$n$ tiles.
\item Every allowed pattern therefore occurs with bounded gaps.
\end{itemize}
\note{$m_n\geq2$ supplies $00,01,11$ inside a parent, while $10$ occurs between parents. Choose $n\geq1$ to avoid the first-level run-length exception.}
\end{frame}

\subsection{Regularity and statistics}

% Manuscript: Lemma 4; Section 3
\begin{frame}{Disagreements shrink by a factor of three}

\[
D_n=\{j:U_n(j)\ne V_n(j)\},\qquad \frac{|D_n|}{H_n}=3^{-n}.
\]
\begin{bookresult}{lem-regular}
The same $Y$ has an aperiodic regular Toeplitz generator.
\end{bookresult}
Each coordinate eventually lies in a common prefix or suffix and acquires a period. Skeleton density is at least $1-3^{-n}$.
\medskip
Extract an essential period structure from the skeletons’ least periods.
\note{Coordinatewise periodicity does not imply a global period. Least skeleton periods form an unbounded divisibility chain; removing repetitions gives an essential period structure. The statements chapter proves the extraction.}
\end{frame}

% Manuscript: Lemma 5
\begin{frame}{One bound determines all statistics}

\begin{bookresult}{lem-unique}
$Y$ has exactly one invariant Borel probability.
\end{bookresult}
\[
\mu(\theta_n=j)=\nu(\theta_n=j)=1/H_n.
\]
A width-$r$ cylinder $C$ has at most $r|D_n|$ bad phases:
\[
|\mu(C)-\nu(C)|\leq r3^{-n}\longrightarrow0.
\]
Fix the window before $n\to\infty$; existence comes from weak limits of empirical orbit probabilities.
\note{On a good phase, the whole phase fibre lies inside or outside the cylinder. Good contributions agree; the difference of bad contributions is at most their common total mass, with no extra factor of two.}
\end{frame}

\section{Splitting}
\subsection{A fixed label}

% Manuscript: Section 4, holonomy equation
\begin{frame}{Each whole block acts trivially}
Any finite state set $Q$, $d=|Q|$, and letter permutations $\pi_0,\pi_1$.
\[
P(a_0\cdots a_{r-1})=\pi_{a_{r-1}}\cdots\pi_{a_0}.
\]
Choose $n\geq\max\{1,d\}$. Each permutation order divides $d!\mid n!$, so
\[
P(U_n)=P(V_{n-1})^{n!}P(U_{n-1})^{2n!}=1,
\]
\[
P(V_n)=1.
\]
Each factor vanishes separately; the permutations need not commute.
\note{The factorial kills the permutation of each entire lower word. The chosen level $n$ may depend on the finite set; the base $Y$ never changes.}
\end{frame}

% Manuscript: Section 4, coboundary and trivialization equations
\begin{frame}{Relabel the state to separate the motion}
$c(y)$ is the accumulated action from the current tile start to just before the origin.
\[
c(Sy)=\pi_{y_0}c(y),\qquad F(y,q)=(y,c(y)^{-1}q).
\]
\[
F(Sy,\pi_{y_0}q)=(Sy,c(y)^{-1}q).
\]
\begin{itemize}
\item Inside a tile: the same left factor cancels.
\item At a boundary: identity whole-tile action permits the reset.
\item Recognizability gives continuity; $F^{-1}(y,q)=(y,c(y)q)$.
\end{itemize}
\note{The boundary check is essential; otherwise the identity holds only inside tiles. The transformed label $c(y)^{-1}q$ is fixed while the base continues shifting.}
\end{frame}

% Manuscript: Section 4; Kang–Jäger Theorem 1.3
\begin{frame}{From a local rule to finite permutations}

\[
Q=\{0,1\}^{\ell-1},\qquad |Q|=2^{\ell-1}.
\]
Given output $a$ and prefix $q$, uniquely solve for the last bit $b$:
\[
h(q_0,\ldots,q_{\ell-2},b)=a,
\]
\[
\pi_a(q)=(q_1,\ldots,q_{\ell-2},b).
\]
The right endpoint gives forward motion; the left gives its inverse. Composing the resulting representation with $F$ gives $\Theta_h$.
\note{At $\ell=2$ the new prefix consists of $b$ alone. Two-sided reconstruction uses finite data for each finite window, giving a homeomorphism on the entire inverse image. Credit Kang--Jäger, Theorem 1.3, for this representation.}
\end{frame}

% Manuscript: Corollary 6
\begin{frame}{Exact components and statistics}

\begin{bookresult}{cor-components}
The full lift has exactly $2^{\ell-1}$ clopen strictly ergodic components. Its invariant-probability simplex has that many ergodic vertices.
\end{bookresult}
\[
\rho=\sum_{q\in Q}t_q\mu_q,\qquad t_q\geq0,\quad\sum_qt_q=1.
\]
Normalize a restriction only at positive weight; zero weight gives the zero restriction. Multiple positive weights yield an invariant slice of intermediate mass.
\note{The topological components and count already have formal proofs. This full written result also includes measure transport and the ergodic-vertex classification. The shared request asks only for those remaining measure clauses.}
\end{frame}

\section{Consequences}

% Manuscript: Section 1; final paragraph of Section 4
\begin{frame}{A negative answer to the question}

\begin{bookresult}{thm-refutation}
Not every strictly ergodic regular Toeplitz binary base admits a minimal complete bipermutative lift.
\end{bookresult}
\begin{itemize}
\item The same $Y$ defeats all lengths and rules.
\item Each component is strictly ergodic; the full lift is neither minimal nor uniquely ergodic.
\item Positive results for other bases remain valid.
\end{itemize}
\note{This separately cites the formal negation of universal existence. It is not expanded into the explicit conjugacy statement. The paper does not classify every base admitting a minimal lift.}
\end{frame}

\appendix

% Manuscript: Bibliography; verification materials
\begin{frame}{Evidence and sources}
Full reading and proofs: Haobo Ma and Wenlin Zhang, version 2.
\medskip
\href{https://doi.org/10.5281/zenodo.23885965}{doi:10.5281/zenodo.23885965}
\begin{itemize}
\item Kang--Jäger, arXiv:2609.01140v1, Question 1.10, Theorems 1.3 and 1.9.
\item Berthé et al., arXiv:2403.17892v3: coboundaries and minimality.
\item Hedlund (1969); Downarowicz (2005): local rules and Toeplitz background.
\item Finite-level inverse continuity and the probability simplex retain complete written proofs.
\end{itemize}
\note{The book provides full references and focused contribution links. Finite checks at initial levels do not replace the infinite-level proof.}
\end{frame}

% Manuscript: Corollary 6, topological clause
\begin{frame}{Counting orbit closures}

\begin{bookresult}{thm-orbit-count}
Every allowed length-$\ell$ rule over the fixed $Y$ has exactly $2^{\ell-1}$ distinct orbit closures.
\end{bookresult}
This is the topological count in the component result. Each orbit closure is one fixed-label slice.
\note{The formal source uses $\ell=k+2$ and prefix set $\operatorname{Fin}(k+1)\to\operatorname{Bool}$, covering all $\ell\geq2$. This theorem does not state measure classification.}
\end{frame}

% Manuscript: Source evidence ledger
\FormalVerification
\end{document}
